\vspace*{-0.25cm}
\section{Conclusion and Future Research}
\label{Sec:concl}

This paper proposed the use of Davio expansions and unbalanced relative phase constructions to optimize relative phase constructions in their usage in LUT-based quantum circuit synthesis. These were meant to improve the $2^n$ proportionality of an earlier, Shannon Decomposition based method. The experimental results showed a clear advantage over this previous method, as well as a significant advantage over the naive method. There remains one interesting exception where the proposed method underperforms the naive method, however, and it relates to a special form of LUT network.

The results of this paper show a clear advantage to using our proposed method, as well as possible further avenues of inquiry for related research. As the experimental results show, there remain many variables to control to optimize in this method, and researching those further could prove fruitful.
%This paper proposed the use of relative phase Boolean functions to optimize LUT-based quantum circuit synthesis while maintaining a minimum of ancillae. While the Shannon decomposition allows us to the cost remains exponential with respect to the number of LUT inputs, limiting the LUT-size to under 6 inputs allows us to see a significant T-count advantage in the proposed method using relative phase Boolean functions.

%There still remain some cases where it underperforms the naive case so further research will be needed there. Furthermore, it remains an open question whether we can get better than exponential cost for the LUTs. As shown in~\cite{bib-lhrs-soeken}, increasing LUT size has possible advantages in gate count and depth. Furthermore, our research has shown that decreasing LUT-size has an advantage in T-count. Existing research in this topic has now uncovered several variables with which to toggle to minimize the T-count, and points to potential further research in the topic.
